% Proved after four complete actual-file audits and exact H2 dependency-discharge refreshes. Coherent band and deck data retained.
\subsection{Bounded-period projective growth with arbitrary field monodromy}

\begin{lemma}[Proved: coherent skew Morita--Frobenius effectivity]
\label{D:skew-Morita-Frobenius-effectivity}
Let $U_V\to U_H$ be the cyclic field cover attached to a
continuous character
$\chi:\pi_1(U_H)\to\operatorname{Gal}(\mathbb F_q/\mathbb F_p)$,
where $q=p^f$, $a=|\operatorname{im}\chi|$, $\nu=f/a$ and
$q_0=p^\nu$. Let $\mathcal G_H$ be an actual lifting gerbe
with band $\mathbb F_q^\times(\chi)$ on $U_H$, and let
$\mathcal G_V$ be its pullback, with the resulting constant
band $\mu_{q-1}$. Write $\mathscr T$ for its coherent deck
pullback, oriented to carry scalar weight one to scalar
weight $q_0$. It covers a deck automorphism $\tau$ of $U_V$,
and its given coherence includes $\mathscr T^a\simeq1$.

Use the deck-equivariant comparable-cycle model $Y$ of
Lemma~\ref{D:equivariant-cover-comparable-cycles}. Its cycles
$C,Z$ and its actual separating pencil $x,y$ are invariant.
Choose rank-$h$ anchors $\mathcal A_j$ of scalar weight $p^j$
for $0\le j<\nu$, with $p\nmid h$, and put
$\mathcal A_\nu=\mathscr T\mathcal A_0$. Define full orders
and comparison bundles as in (D.MF.1), using this terminal
anchor. Thus
\[
 P_\nu=\tau^*P_0,\qquad
 K_j=\operatorname{full}j_*\mathcal Hom
              (\mathcal A_{j+1},F^*\mathcal A_j),\qquad
 R_j=\operatorname{full}j_*\mathcal Hom
              (F^*\mathcal A_j,\mathcal A_{j+1})(-pC).
                                                        \tag{D.SM.1}
\]
Here all scalar-weight-zero sheaves descend before their
full pullback, and $F$ is absolute $p$-Frobenius.

Take $b,N,L$ exactly as in (D.MF.2). On $LC_T$ over any
coefficient algebra $T$, take the same finite data and
actual action, balance, intertwining and two right-inverse
equations (D.MF.3), (D.MF.6), for $0\le j<\nu$, except that
\[
             E_\nu=\tau^*E_0,\qquad
 \mu_{\nu-1}:F^*E_{\nu-1}\otimes K_{\nu-1}
                                     \longrightarrow\tau^*E_0 .
                                                        \tag{D.SM.2}
\]
These data have actual formal realizations with their
specified $NC_T$ marking and with genuine skew Morita
unit maps on the puncture. Two realizations with the same
marked $NC_T$ data give isomorphic punctured unit objects.
For finite-type Noetherian or finite local Artinian
coefficient bases, the ordinary bundles and their structure
maps algebraize over the same completed tensor base as in
Theorem~\ref{D:finite-Morita-Frobenius-object-effectivity}.

For actual families with every $E_j(-Z)$ generated and
every determinant fixed formally, the analogue of
Proposition~\ref{D:generated-Morita-unit-Isom-openness}
holds: the fixed-object meromorphic unit-Isom functor has
finite presentation and is etale, and its class image is
open and closed. The punctured objects are exactly the
twisted projective-semilinear frame objects on the specified
coefficient gerbe. No gerbe extension to $Y$, root-lattice
flatness, averaging by the order $a$, or order-$a$ matrix
equation is assumed.
\end{lemma}

\begin{proof}
We check the changes to the already proved simultaneous
construction, including its categorical comparison.
The actual banded gerbe is pulled back from $U_H$, so
its deck functors and all their coherence maps are given
by descent; they are not selected from a geometric Brauer
class alone. Pullback commutes with absolute Frobenius.
The terminal scalar weight is $p^\nu=q_0$, so every Hom
in (D.SM.1) has weight zero, as required.
In particular
\[
 \mathcal Hom(\mathscr T\mathcal A_0,\mathscr T\mathcal F_0)
       =\tau^*\mathcal Hom(\mathcal A_0,\mathcal F_0)
                                                        \tag{D.SM.3}
\]
after descent to $U_V$. Full pullback commutes with the
automorphism of the same regular model. This proves the
typing of the terminal order and module; no extra
Frobenius operation occurs in $\tau^*$.

The source-only Frobenius comparison (D.MF.4) and the
actual Morita conductor hold for every consecutive pair
of the anchors in (D.SM.1), including the terminal pair.
The opposite $-pC$ twist gives the same conductor
$bC=(p+1)C$. The global tensors for $x^b,y^b$ are actual
Noether-factorization tensors, exactly as in (D.MF.5).
The deck group fixes $x,y$, since these functions come
from the fixed base surface. There is thus no change to
the two unlocalized inverse equations on $LC_T$.

Choose a finite affine formal trivializing cover closed
under $\tau$; its size is not a parameter count. The
finite-cycle frames, affine nilpotent-thickening lifts,
and arbitrary-coefficient Cartier regularity used in the
proof of Theorem~\ref{D:finite-Morita-Frobenius-object-effectivity}
apply on this refinement. Use the corrected actions at
the initial stage and their deck pullbacks at the terminal
stage, rather than independently correcting a duplicate
terminal action. The two scaled-separability exponents
are still $2$ and $2p$. Action Newton retains $L-2$,
the two commuting averages retain
$M=L-2p-4=N+b$, and the exact inverse normalization
changes the opposite map only modulo $NC_T$. These are
the same actual algebra identities, with the last target
replaced by its deck pullback. The induced balanced
source and target have equal ordinary rank $hr$ on
$U_V$, so the right inverse is again an actual Morita
isomorphism.

For $j<\nu-1$ the overlap operator is (D.MF.9).
At the last stage it is, with the corresponding affine
indices understood,
\[
 \mathcal T(g)_0=(\tau^{-1})^*
 \left(\mu_{\nu-1,i}
       (F^*g_{\nu-1}\otimes1)\mu_{\nu-1,j}^{-1}\right).
                                                        \tag{D.SM.4}
\]
It preserves the exact order-intertwining relation and
is functorial for products and inverses. Since $\tau$
preserves the ideal of $C$, every error of order $u$ is
sent to order at least $pu-b$, including at (D.SM.4).
The initial overlap defect is still $M-b=N$; this
uses the unchanged forward map, even though the opposite
map was corrected only to precision $N$.
The same strict inequalities in (D.MF.2) give convergence,
unique overlap limits, actual cocycles, and comparison
of two realizations with identical $NC_T$ data. The
proof uses no averaging over the deck group, so it
also covers $p\mid a$.
All regular tensor equations glue by Cartier
torsion-freeness. Formal effectivity over arbitrary
coefficient algebras and Noetherian Grothendieck existence
therefore have exactly their previous distinct scopes.
Azumaya Morita is used only on $U_V$.

For clarity, the resulting twisted diagrams are
\[
 F^*\mathcal F_j\xrightarrow{\sim}\mathcal F_{j+1}
       \ (j<\nu-1),\qquad
 F^*\mathcal F_{\nu-1}\xrightarrow{\sim}
                                      \mathscr T\mathcal F_0 .
                                                        \tag{D.SM.5}
\]
Their composite is a map
$F_{q_0}^*\mathcal F_0\to\mathscr T\mathcal F_0$.
For weight-zero sheaves this is
$F_{q_0}^*L\to\tau^*L$, equivalently
$(\tau^{-1})^*F_{q_0}^*L\to L$. Thus the generator in
the ordinary skew-frame convention of
Lemma~\ref{D:twisted-deck-Frobenius-semilinear-model}
is $\tau^{-1}$ with our terminal-target convention.
Iterate it $a$ times, using the \emph{given} identification
$\mathscr T^a\simeq1$. This defines an $F_q$-unit map.
Its finite etale $\mathbb F_q$-frame torsor inherits the
semilinear deck transport, and the $a$-fold deck identity
holds on its fixed frames. Conversely finite semilinear
frames on the actual lifting gerbe give (D.SM.5).
This is precisely the etale-atlas categorical argument
at the end of Lemma~\ref{D:compressed-skew-Frobenius-unit-jets}
and Lemma~\ref{D:twisted-deck-Frobenius-semilinear-model};
that argument uses the punctured coefficient gerbe and
its coherent descent, not its extension over a root
stack or any tameness assumption. Morphisms and
coefficient base changes are respected.
The $a$-fold iteration is used for this identification
only, never as an expanded polynomial equation.
An arbitrary pointed lift of the deck generator need
not have order $a$.

Finally consider comparison of actual families.
The global averaging operators and the finite $S_x,S_y$
equations in
Proposition~\ref{D:generated-Morita-unit-Isom-openness}
still exist. At the terminal stage their target is
$\tau^*$ of the initial Hom bundle, or equivalently
one applies $(\tau^{-1})^*$ afterward. Both pencil
elements and $C$ are invariant, so their loss remains
$d=2p+b$; a meromorphic Hom lattice shifted by $\ell C$
has loss $d+(p-1)\ell$ as before. Consequently their
intrinsic finite fixed-point equations, paired inverse
equations and proper Hom finite-presentation argument
are unchanged.

The target-Hartogs and determinant-adjugate comparison
uses only the generated source, fixed determinant,
and canonical fixed target at each of the $\nu$ stages.
It therefore supplies the same reference-dependent
pole bound for all specified Artinian isomorphisms,
independent of nilpotence. The bound is qualitative,
not an input to the parameter dimension. On the
coefficient gerbe finite etale frames and specified
isomorphisms lift uniquely across nilpotent thickenings.
Their deck compatibility follows from uniqueness.
Thus the represented comparison is formally etale
and of finite presentation. Its image is open, and
its complement is the union of the other open class
images. This proves the stated clopen comparison.
\end{proof}

\begin{theorem}[Proved: bounded-period hereditary projective semilinear growth]
\label{D:bounded-period-semilinear-hereditary-projective-growth}
Let $(R,\mathfrak m,k)$ be a fixed complete excellent
normal local surface over an algebraically closed
coefficient and residue field of characteristic $p>0$,
and put $\Gamma=\pi_1(\operatorname{Spec}R\setminus\{\mathfrak m\})$.
Fix $B\ge1$.
Then there is $A_{R,B}^{\rm ps}$ such that, for every
open subgroup $H\subset\Gamma$ of index $t$, every
$q=p^f$, every $r\ge1$, and every fixed field character
$\chi:H\to\operatorname{Gal}(\mathbb F_q/\mathbb F_p)$,
the homomorphisms
\[
 H\longrightarrow\operatorname{PGL}_r(\mathbb F_q)
                 \rtimes\operatorname{Gal}(\mathbb F_q/\mathbb F_p)
       \quad\text{with field projection }\chi
                                                        \tag{D.SM.6}
\]
have at most $q^{A_{R,B}^{\rm ps}tr^2}$ projective
conjugacy classes in the following specified range:
the geometric Brauer obstruction after restriction
to $K=\ker\chi$ has order at most $B$.
There is no tameness assumption on either cover and
no bound on the field-image order. The range includes
every representation (D.SM.6) if $\gcd(r,q-1)\le B$.
Unbounded geometric period is not asserted.
\end{theorem}

\begin{proof}
The case $r=1$ is immediate, and $q=2$ is the ordinary
hereditary count. Otherwise put $n=q-1$,
$a=|\operatorname{im}\chi|$, $\nu=f/a$, and $q_0=p^\nu$.
The field cover $U_V\to U_H$ has degree $a$ and total
degree $ta$ over the original puncture.

The exact sequence
$1\to\mathbb F_q^\times\to\operatorname{GL}_r(\mathbb F_q)
\rtimes\operatorname{Gal}(\mathbb F_q/\mathbb F_p)
\to\operatorname{PGL}_r(\mathbb F_q)
\rtimes\operatorname{Gal}(\mathbb F_q/\mathbb F_p)\to1$
associates to (D.SM.6) a class with its actual coefficient
action in $H^2(H,\mathbb F_q^\times(\chi))$.
The group-to-etale comparison is injective.
Lemma~\ref{D:arbitrary-field-character-H2} therefore
bounds the number of realized actual coefficient classes by
\[
                 |H^2(U_H,\mathbb F_q^\times(\chi))|
                                  \le q^{c_Rt}.        \tag{D.SM.7}
\]
Fix such a class and its actual lifting gerbe.
It is not replaced by a gerbe for its geometric image.

On the field cover its band becomes $\mu_n$.
If its geometric Brauer period is $M\le B$, the rank
$h=M^2$ anchor construction used in
Theorem~\ref{D:bounded-period-hereditary-projective-growth}
gives rank-$h$ anchors for every invertible scalar
weight. Choose weights $p^j$, $0\le j<\nu$, and take
the terminal anchor to be $\mathscr T\mathcal A_0$.
Only the existence on the puncture is used.
In particular $p\nmid h$. The anchors need not have
a finite-field unit structure or a separately chosen
deck linearization.

Use the equivariant comparable-cycle model on the
degree-$ta$ cover. On the same original model the
canonical full orders $P_j$, comparison bundles
$K_j,R_j$, and modules $E_j$ of geometric inputs have
the actual actions and conductors of (D.SM.1)--(D.SM.2).
Their ranks are $h^2,h^2,hr$, respectively.
The full-bundle degree and generation theorem gives
\[
 E_j(-Z)\ \text{generated},\qquad
 0\le c_1(E_j)\cdot Z=O_{R,B}(ta r),\qquad
 h^0(\mathcal O_{\ell C})+h^1(\mathcal O_{\ell C})
                                      =O_{R,\ell}(ta).
                                                        \tag{D.SM.8}
\]
Fixed order and comparison terms have degree
$O_{R,B}(ta)$. These bounds concern the ordinary
full bundles on this model; there is no root-node
rectangle or intermediate-lattice generation argument.

\emph{Determinants before passage to the field cover.}
For a nonempty fixed coefficient/rank class choose
one twisted skew-unit reference $\mathcal F^*$.
With the same anchors the punctured ratios satisfy
\[
 \frac{\det E_j}{\det E_j^*}
   =\left(\frac{\det\mathcal F_j}
                    {\det\mathcal F_j^*}\right)^h .
                                                        \tag{D.SM.9}
\]
The ratios in parentheses have scalar weight zero.
Their $\nu$ Frobenius arrows and terminal coherent
deck arrow give one rank-one semilinear local system
on $U_H$. Its class lies in
$H^1(H,\mathbb F_q^\times(\chi))$.
Lemma~\ref{D:arbitrary-field-character-scalar-H1}
bounds all initial choices by $q^{c_Rt}$.
Each of the remaining phases is determined by one
absolute Frobenius pullback. No enumeration of
ordinary characters on the larger degree-$ta$ cover
is substituted for this crossed-character bound.

There are $O_R(ta)$ exceptional components. The
nonnegative multidegrees of each canonical full $E_j$
have sum $O_{R,B}(ta r)$ by (D.SM.8) and $C\le Z$.
For a fixed punctured line its extension is uniquely
determined by this vector, by negative definiteness.
The ordinary determinant enumeration thus gives at most
\[
 q^{c_Rt}
    \exp\!\left(c_{R,B}\nu ta\{\log(r+1)+1\}\right)
       =q^{c_Rt}
    \exp\!\left(c_{R,B}ft\{\log(r+1)+1\}\right)          \tag{D.SM.10}
\]
lists of actual global determinant lines.
Retain only lists occurring for a canonical geometric
input, so the corresponding determinant lines
$L_j=\det(E_j(-Z))$ are actually generated.
The terminal list is $\tau^*L_0$ and is not chosen
independently. The fixed anchor comparison contributes
no further determinant-local-system choice.

\emph{The finite polynomial model.}
Choose $N,L$ as in the preceding lemma, independent
of $a,t,r,f$ for the fixed original geometry and $B$,
and put $D=LC$. For each retained determinant list use
the actual quotient presentations
\[
 0\longrightarrow L_j^{-1}|_D\longrightarrow
       \mathcal O_D^{hr+1}\longrightarrow V_j\longrightarrow0,
       \qquad E_j=V_j(Z).
                                                        \tag{D.SM.11}
\]
The open subbundle condition gives locally split
sequences over every coefficient algebra.
Introduce one actual order action per stage and the
balanced forward/opposite Morita maps, and impose
their action, associativity, balance, and two scaled
right-inverse equations. All maps have the precise
typing in the preceding lemma.

The full fixed-reference quotient/Hom proof of
Theorem~\ref{D:bounded-period-hereditary-projective-growth}
applies with cover degree $ta$. Its top--top, mixed,
and bottom--bottom blocks have multiplicities
$O_B(r^2),O_B(r),O_B(1)$ and corresponding line
cohomology budgets $O_{R,B}(ta),O_{R,B}(ta r),
O_{R,B}(ta r)$. Fixed reference tensor words have
bounded rank and bounded internal width in $B$.
Only two variable object ends occur, including in
associativity and composition identities. Consequently
each stage has $O_{R,B}(ta r^2)$ object and equation
coordinates and bounded polynomial degrees.

At the terminal stage the target quotient, determinant,
order and reference maps are their actual $\tau$
pullbacks. The parameter coefficients are the same
as at stage zero; the only nonlinear coefficient
operation is the adjacent $p$-Frobenius.
The pulled-back fixed line is generated, with the
same cohomology bounds because $C,Z$ are invariant.
One computes its Hom complexes on any fixed two-affine
cover, treating the pulled-back lines and maps as fixed
reference data. A deck-stable trivializing refinement
is used only for qualitative effectivity and is not
multiplied into this coordinate count.
The all-ring finite Hom and boundary-witness construction
therefore yields one polynomial zero scheme and open
subbundle locus with
\[
 N_{\rm var}\le c_{R,B}\nu ta r^2=c_{R,B}ft r^2,
 \qquad
 \deg(\text{each equation})\le D_{R,B}.
                                                        \tag{D.SM.12}
\]
Each equation uses a bounded number of adjacent phases.
Neither a degree-$q$ polynomial nor an expanded
$a$-fold deck word appears.

\emph{Effectivity, class comparison, and coverage.}
Lemma~\ref{D:skew-Morita-Frobenius-effectivity} turns
every point of this finite model into an actual formal
skew unit object with the prescribed $NC$ marking.
Noetherian completed-base existence algebraizes it
where required. Its retained quotient on $Z\subset NC$
lifts globally, and its determinant remains the fixed
line by the weak-cycle generator/Picard lemma.
This holds for the whole parameter family, not only
for its initially canonical geometric points.

Every original projective-semilinear representation
occurs. Its canonical full input supplies one retained
determinant list and the finite data (D.SM.11).
That genuine input and the constructed realization
have the same $NC$ marking. The comparison in the
preceding lemma identifies their punctured skew
unit objects, so finite reconstruction does not
change the represented class.

For a fixed represented class, use its canonical
full coarse modules as fixed comparison targets.
The generated-source Hartogs extension, fixed
determinant ratio and adjugate bound give a uniform
pole bound for every specified Artinian comparison,
independent of nilpotence. The bound may depend on
this class and rank, which is harmless for qualitative
finite presentation. The preceding lemma's global
finite $S_x,S_y$ equations and paired inverses
represent these comparisons; unique nilpotent
lifting of finite etale frames preserves the actual
coefficient-band and deck coherence.
The fixed-object Isom map is therefore etale of
finite presentation, with clopen class image.
Each irreducible component of the parameter open
contains at most one such class.

The component bound applied to (D.SM.12) is
$(N_{\rm var}+1)\max\{2,D_{R,B}\}^{N_{\rm var}}
\le p^{c_{R,B}ft r^2}$.
Multiplying by the coefficient classes (D.SM.7)
and determinant lists (D.SM.10) gives
$q^{A_{R,B}^{\rm ps}tr^2}$.
Fixed frames on the actual lifting gerbe project
onto the projective-semilinear representation classes
with that coefficient obstruction. Extra choices
of a twisted lift only increase the set being counted,
so they do not invalidate this upper bound.

Finally restriction to $K$ is an ordinary projective
representation into $\operatorname{PGL}_r(\mathbb F_q)$.
Its scalar multiplier is killed by both $r$ (take
determinants of lifts) and $q-1$ (the coefficient
exponent). Its geometric period therefore divides
$\gcd(r,q-1)$, proving the sufficient condition.
The construction includes $a=1$ and $p\mid a$.
\end{proof}

\begin{remark}[Exact remaining condition]
The period bound is imposed on the ordinary geometric
Brauer class on the field cover. It is not a claimed
ordinary Brauer class for the nonconstant band on $U_H$,
nor a bound on the full coefficient obstruction.
The twisted $H^2$ input is supplied by
Lemma~\ref{D:arbitrary-field-character-H2} in its proved
arbitrary-cover scope. Neither result here asserts
unbounded-period projective growth.
\end{remark}
